\section{Planning}
\label{sec:planning}

This section outlines the planning of the research activities described in Section~\ref{sec:research-design}. The study is planned over twenty-one weeks and is structured according to the main activities of the design science process, from problem identification and objective definition to artifact development, demonstration, evaluation, and communication.

Although the activities are shown in a largely sequential order, the research process is iterative. Insights from demonstration and evaluation may lead to refinements of earlier design decisions, for example in the structured representation of user stories, the linking rules, or the coverage assessment procedure. Writing is treated as a continuous activity throughout the project to ensure timely documentation of design decisions, intermediate results, and final findings.

Figure~\ref{fig:planning-gantt} summarizes the planning. The first phase focuses on problem clarification, solution objectives, and preparation of the required data and materials. The middle phase is dedicated to the design and implementation of the traceability method and prototype. The final phase focuses on demonstration, evaluation, and consolidation of the results in the thesis.

\begin{figure}[h!tb]
    \centering
    \resizebox{\textwidth}{!}{
        \begin{ganttchart}[
            hgrid,
            vgrid
        ]{1}{21}
            \gantttitle{Weeks}{21} \\
            \gantttitlelist{1,...,21}{1} \\

            \ganttbar{DSRM-A1: Problem identification and motivation}{1}{3} \\
            \ganttbar{DSRM-A2: Definition of solution objectives}{3}{5} \\
            \ganttbar{Data and materials preparation (D1--D4)}{4}{6} \\
            \ganttbar{DSRM-A3: Design and development}{6}{13} \\
            \ganttmilestone{Prototype implementation (A3)}{13} \\
            \ganttbar{DSRM-A4: Demonstration}{14}{16} \\
            \ganttbar{DSRM-A5: Evaluation}{17}{19} \\
            \ganttbar{DSRM-A6: Communication (thesis writing)}{6}{21} \\
            \ganttmilestone{Thesis submission}{21}
        \end{ganttchart}
    }
    \caption{Planning of the research activities over twenty-one weeks, aligned with the main design science phases.}
    \label{fig:planning-gantt}
\end{figure}

% \section{Planning}
% \label{sec:planning}

% This section outlines the planning of the research activities described in Section~\ref{subsec:research-method}. The planning follows the Design Science Research Methodology (DSRM) activities (DSRM-A1--A6) and indicates how these activities are distributed over time.

% The research is planned over twenty-one weeks. Although the activities are presented sequentially, the research process is iterative: insights from demonstration and evaluation may lead to refinements of earlier design decisions (e.g., adjustments to the structured story representation or tracing approach). Communication (writing) is continuous to ensure timely documentation of design decisions and results.

% \begin{figure}[h!tb]
%     \centering
%     \resizebox{\textwidth}{!}{
%         \begin{ganttchart}[
%             hgrid,
%             vgrid
%         ]{1}{21}
%             \gantttitle{Weeks}{21} \\
%             \gantttitlelist{1,...,21}{1} \\

%             \ganttbar{DSRM-A1: Problem identification and motivation}{1}{3} \\
%             \ganttbar{DSRM-A2: Definition of solution objectives}{3}{5} \\
%             \ganttbar{Data and materials preparation (D1--D4)}{4}{6} \\
%             \ganttbar{DSRM-A3: Design and development}{6}{13} \\
%             \ganttmilestone{Prototype implementation (A4)}{13} \\
%             \ganttbar{DSRM-A4: Demonstration}{14}{16} \\
%             \ganttbar{DSRM-A5: Evaluation}{17}{19} \\
%             \ganttbar{DSRM-A6: Communication (thesis writing)}{6}{21} \\
%             \ganttmilestone{Thesis submission}{21}
%         \end{ganttchart}
%     }
%     \caption{Planning of the research activities aligned with the DSRM phases.}
%     \label{fig:planning-gantt}
% \end{figure}
